% Source: https://artofproblemsolving.com/community/c3h2348985_yeetopedias_problem_thread
% Stable identity: aops-topic-2348985-q01
Write $N=10k+1$. The second condition says $10k+1\equiv3\pmod{11}$. Since $10\equiv-1\pmod{11}$, this is $-k\equiv2$, that is $k\equiv-2\equiv9\pmod{11}$. So $k=9+11t$ and
\[N=10(9+11t)+1=91+110t,\qquad t=0,1,2,\dots\]
Conversely every such number satisfies both conditions: $91=9\cdot10+1=8\cdot11+3$ and $110$ is a multiple of both $10$ and $11$.

\textbf{(a)} Since $91=7\cdot12+7$ and $110=9\cdot12+2$,
\[N\equiv7+2t\pmod{12}.\]
For $t=0,1,2,3,4,5$ this gives the remainders $7,9,11,1,3,5$ (for example $N=91,201,311,421,531,641$), and then they repeat. So the remainder can be any odd number from $0$ to $11$: $\boxed{1,3,5,7,9,11}$. (It must be odd, since $N$ is odd.)

\textbf{(b)} The smallest such $N$ is $91=7\cdot12+7$, with remainder $\boxed{7}$.
